% Proposed experiments only; TBD cells are not measured outcomes.
\begingroup
\newcommand{\pending}{\textit{TBD}}

\section{Planned experiments: resolution and efficiency}
\label{sec:planned-outcomes}
These tables specify the experiments needed to extend the paper's argument.
\textit{TBD} denotes an unmeasured result, not a projected score. Sample sizes
are proposed allocations. All new solver experiments use the authorized
DeepSeek V4 Pro OpenRouter route; new embeddings use the pinned Qwen3-Embedding-8B
space. Each submission is verified immediately with its benchmark's verifier;
independent semantic trajectory reviews occur at most ten minutes apart.
Intervals resample paired targets, retaining each target's repetitions together.
Infrastructure failures remain unassessed; paired effects state the number of
jointly assessed targets. USD denotes metered API charges; local CPU/GPU time
is measured separately.

\begin{table}[htbp]
\centering
\caption{\textbf{Does the complete method improve repair success?}
Planned matched comparison on the four fixed manifests. Each cell will report
resolved/planned and assessed/planned for one full first pass. Every arm uses
the same agent, model, tools, task order, and total per-task budget, including
memory composition and execution. Flat and Graph share the frozen bank built
from source273 plus the admitted expansion in Table~\ref{tab:planned-coverage}.
Flat returns complete records by FAISS; Graph uses the selected ranking,
relations, composition, and feedback. The final row reports their resolution
difference on jointly assessed pairs, its 95\% interval, and paired coverage.
Table~\ref{tab:planned-graph} separates the components of this system comparison.}
\label{tab:planned-fourbench}
\small
\begin{tabular}{lcccc}
\toprule
Method & Verified & Related-Lite & LoLBench & DeepSWE \\
Planned tasks & 150 & 99 & 100 & 113 \\
\midrule
No memory & \pending & \pending & \pending & \pending \\
FAISS Flat & \pending & \pending & \pending & \pending \\
ContextGraph & \pending & \pending & \pending & \pending \\
\midrule
Graph $-$ Flat (pp, CI) & \pending & \pending & \pending & \pending \\
\bottomrule
\end{tabular}
\end{table}

\begin{table}[htbp]
\centering
\caption{\textbf{Does memory reach useful solutions with less computation?}
Extend the same frozen comparison to three real passes, retrying each arm's
own validly unresolved tasks. Produce this panel separately for each benchmark.
Freeze memory across passes and retain each arm's own prior-attempt feedback
under the same policy. Report cumulative resolved/planned, all-attempt tokens
and USD per resolved task after three passes, and active seconds to resolve 10\% of the
fixed manifest. Active time includes model, tool, and verification execution;
worker allocation and task order are matched. A threshold never reached is
reported as such. The 10\% threshold is rounded up to a whole task. Count
target-time composition, failed attempts, and verification in active time;
zero resolutions give undefined cost per solve. Add the one-time build cost
from Table~\ref{tab:planned-amortization} to the three-pass API total. This tests the early-resolution
pattern in Table~\ref{tab:swectx-stream} using measured time and cost.}
\label{tab:planned-efficiency}
\small
\begin{tabular}{lcccccc}
\toprule
Method & Pass@1 & Pass@2 & Pass@3 & Tokens/solve & USD/solve & Time to 10\% \\
\midrule
No memory & \pending & \pending & \pending & \pending & \pending & \pending \\
FAISS Flat & \pending & \pending & \pending & \pending & \pending & \pending \\
ContextGraph & \pending & \pending & \pending & \pending & \pending & \pending \\
\bottomrule
\end{tabular}
\end{table}

\begin{table}[htbp]
\centering
\caption{\textbf{Order the work around the scientific argument.}
The experiment sequence uses development tasks to determine useful content and
fix the method before the four-manifest comparison. Each later stage answers
a distinct question; numerical gains are measured rather than assumed.}
\label{tab:planned-sequence}
\small
\begin{tabular}{lp{0.40\linewidth}p{0.40\linewidth}}
\toprule
Stage & Question & Experiment and evidence \\
\midrule
1 & What should memory communicate? & Content and source-utility interventions
in Tables~\ref{tab:planned-content}--\ref{tab:planned-utility}. \\
2 & What must the graph preserve? & Relation and composition ablations in
Table~\ref{tab:planned-graph}. \\
3 & Can retrieval find useful conditions? & Ranking and source coverage in
Tables~\ref{tab:planned-retrieval}--\ref{tab:planned-pool}. \\
4 & Does the fixed method help in practice? & Four-benchmark resolution,
retry efficiency, and amortized cost in Tables~\ref{tab:planned-fourbench},
\ref{tab:planned-efficiency}, and~\ref{tab:planned-amortization}. \\
\bottomrule
\end{tabular}
\end{table}

\clearpage
\section{Planned experiments: which memory is effective?}
\label{sec:planned-content}
Use 24 distinct development targets from the existing dev70 partition, spanning
at least four repositories and four operation families. Select task/source
pairs from public issues and source evidence before observing the new repairs.
Use three fresh repetitions per arm, with common target and source IDs, and
counterbalance arm order within targets. The proposed allocation is 72 attempts
per arm. Fix a common target-level evaluation suite before candidate outcomes:
each target has applicable joint checks and neighboring behaviors that pass
on its original base. These tests define the outcomes across content, source,
graph, and retrieval interventions. A regression is an attempt failing at
least one of these originally passing preservation checks.

\begin{table}[htbp]
\centering
\caption{\textbf{Which representation transfers the missing requirement?}
Each target uses the same single historical source and the same solve budget.
The rows vary only the delivered content: requirement $r$, historical patch
$H$, witness $W=(I,A,O)$, and source executions $E^{-},E^{+}$.
All rows receive the same execution tool and current-task instructions.
Official is benchmark resolution; Joint requires all predeclared combined
current/historical checks; Preserve requires the neighboring input checks;
Combined requires all three on the same repair in this fixed suite. Additional joint
and neighboring evaluation inputs and observers are fixed before these runs
and withheld from the solver; delivered source witnesses remain visible.
Report mean task success across repetitions and paired intervals. The
$r+W$ versus $r+H$ contrast compares executable examples with a patch reference;
the last two rows isolate supplying the source execution observations.}
\label{tab:planned-content}
\small
\begin{tabular}{lcccc}
\toprule
Delivered memory & Official $\uparrow$ & Joint $\uparrow$ & Preserve $\uparrow$ & Combined $\uparrow$ \\
\midrule
None & \pending & \pending & \pending & \pending \\
Requirement $r$ & \pending & \pending & \pending & \pending \\
Historical implementation $H$ & \pending & \pending & \pending & \pending \\
$r+H$ & \pending & \pending & \pending & \pending \\
$r+W$ & \pending & \pending & \pending & \pending \\
$r+W+E^{-}+E^{+}$ & \pending & \pending & \pending & \pending \\
\bottomrule
\end{tabular}
\end{table}

\begin{table}[htbp]
\centering
\caption{\textbf{What makes one source experience useful?}
On a 12-target subset chosen before candidate outcomes, form a pool of four
distinct sources per target using dense, lexical, operation-connected, and
same-repository random selection, filling duplicates from the same policy.
Deliver each source individually in the complete-witness format for three
repetitions: 144 source-conditioned attempts, plus the matched no-memory
controls from Table~\ref{tab:planned-content}, evaluated with the same fixed suite.
First separate sources without an established operation connection. Classify
connected sources by compatibility with the current requirement, then by whether
the behavior holds on the original workflow. Report source-pair
counts, change in Combined success, and new failures of the originally passing
preservation checks. These per-source effects supply development
utility labels for Table~\ref{tab:planned-retrieval}.}
\label{tab:planned-utility}
\small
\begin{tabular}{p{0.44\linewidth}ccc}
\toprule
Source/current relation & Source pairs & $\Delta$ Combined & $\Delta$ regressions \\
\midrule
Adds an omitted requirement & \pending & \pending & \pending \\
Preserves already-satisfied behavior & \pending & \pending & \pending \\
Conflicts with the current requirement & \pending & \pending & \pending \\
No established operation connection & \pending & \pending & \pending \\
\bottomrule
\end{tabular}
\end{table}

\clearpage
\section{Planned experiments: graph relations and retrieval}
\label{sec:planned-graph}

\begin{table}[htbp]
\centering
\caption{\textbf{Which graph operation improves the submitted repair?}
Use the same 24 development targets and three repetitions, fixing the FAISS
top-three source IDs, complete source content, memory token allowance, current
feedback, and callable tool. A delivers the source witnesses. B asks a writer
to generate joint checks from the same text records and public example;
C uses explicit input and operation bindings with the supported adapters.
D and E mask one binding family while preserving source text and values;
the same bounded writer reconstructs the missing bindings before the same
adapter runs. Thus they test supplying versus reconstructing a relation.
B--E request the same number of checks. F adds their actual failed observations
to C. Compare B--A for generated checks, C--B for graph-bound construction,
C--D/E for supplied relations, and F--C for feedback. Match the full generation,
execution, and solver budget. Failed construction remains in the target
denominator; report valid joint-check delivery as well as repair outcomes.}
\label{tab:planned-graph}
\small
\begin{tabular}{p{0.38\linewidth}cccc}
\toprule
Variant & Joint delivery & Joint & Combined & Regressions \\
\midrule
A. Flat witnesses; no composition & \pending & \pending & \pending & \pending \\
B. A + text-based joint checks & \pending & \pending & \pending & \pending \\
C. A + graph-bound composition & \pending & \pending & \pending & \pending \\
D. Reconstruct C's input-role bindings & \pending & \pending & \pending & \pending \\
E. Reconstruct C's operation bindings & \pending & \pending & \pending & \pending \\
F. C + executed-condition feedback & \pending & \pending & \pending & \pending \\
\bottomrule
\end{tabular}
\end{table}

\begin{table}[htbp]
\centering
\caption{\textbf{Can a retriever find behaviorally useful memory?}
Panel A ranks the four-source pools from Table~\ref{tab:planned-utility};
panel B ranks the entire frozen source bank on all 24 development targets.
For B, measure standalone source utility on the union of all policies' top-three
sources, using the same intervention as Table~\ref{tab:planned-utility}.
Freeze rankings and the evaluation suite before obtaining these utility labels.
Hold the payload format and maximum $k=3$ fixed. Utility@3 is the sum of
standalone $\Delta$ Combined over delivered sources divided by three; empty
slots contribute zero, the no-memory utility. Report uncertainty across targets.
A full graph policy is then fixed on development
data for Table~\ref{tab:planned-fourbench}. As an end-to-end development check,
run each top-three packet with three repetitions and the
same composer; the Combined column measures the packet's actual combined
effect. Composable@3 is the fraction of targets with at least one retrieved
condition that retains the current and historical triggers and executes with
its observer. A target with no delivered condition remains in this denominator.
The operation-reranking row is a proposed policy; the current implementation
retains FAISS order and supplies operation links without a hard veto.}
\label{tab:planned-retrieval}
\small
\begin{tabular}{p{0.40\linewidth}cccc}
\toprule
Candidate ranking & Utility@3 & Composable@3 & Combined & Retrieval ms \\
\midrule
\multicolumn{5}{l}{A. Controlled four-source pools; 12 development targets} \\
\midrule
FAISS dense similarity & \pending & \pending & \pending & \pending \\
BM25 & \pending & \pending & \pending & \pending \\
Dense + BM25 reciprocal-rank fusion & \pending & \pending & \pending & \pending \\
Dense + soft operation reranking & \pending & \pending & \pending & \pending \\
Dense with a required operation match & \pending & \pending & \pending & \pending \\
\midrule
\multicolumn{5}{l}{B. Full source bank; 24 development targets} \\
\midrule
FAISS dense similarity & \pending & \pending & \pending & \pending \\
BM25 & \pending & \pending & \pending & \pending \\
Dense + BM25 reciprocal-rank fusion & \pending & \pending & \pending & \pending \\
Dense + soft operation reranking & \pending & \pending & \pending & \pending \\
Dense with a required operation match & \pending & \pending & \pending & \pending \\
\bottomrule
\end{tabular}
\end{table}

\clearpage
\section{Planned experiments: coverage, scale, and reuse}
\label{sec:planned-coverage}

\begin{table}[htbp]
\centering
\caption{\textbf{Can the bank supply usable experience across the four benchmarks?}
The source273 column is an existing measured same-repository witness-coverage
count, meaning at least one witnessed source from the repository. Expand the
bank from official training data and pre-base commits, and develop the operation
adapters needed beyond Django before the method freeze. New witnesses count
source false-to-true contrasts. Expanded coverage and valid joint delivery use
all planned targets as their denominator; the latter also requires both triggers
and a working observer. An unsupported composition falls back to the stored
witness and is counted separately. Freeze this delivery policy with the common
bank before the four-benchmark comparison.}
\label{tab:planned-coverage}
\small
\begin{tabular}{lcccc}
\toprule
Benchmark & Source273 coverage & New witnesses & Expanded coverage & Joint delivery \\
\midrule
Verified150 & 133/150 & \pending & \pending & \pending \\
Related-Lite99 & 91/99 & \pending & \pending & \pending \\
LoLBench100 & 2/100 & \pending & \pending & \pending \\
DeepSWE113 & 0/113 & \pending & \pending & \pending \\
\bottomrule
\end{tabular}
\end{table}

\begin{table}[htbp]
\centering
\caption{\textbf{Does scale help through useful-source coverage or ranking?}
Use the 12-target source-utility study with a fixed delivered $k=3$.
First add 0, 4, or 16 same-repository distractors per target while retaining
its original four-source pool; measure selection stability and Combined
success for FAISS and the chosen graph policy. Then replace the source with
the highest measured standalone utility with an additional distractor,
keeping the pool size fixed. Freeze these manipulations before new packet
outcomes and run three repetitions per cell. Distractors have complete witnesses
for other operations, selected from public source bindings. This development
experiment uses earlier standalone utility labels and fresh packet trials.
Retention is reported separately for Flat and Graph (F/G); choosing a source
by measured utility does not imply its effect is positive.}
\label{tab:planned-pool}
\small
\begin{tabular}{lcccc}
\toprule
Pool condition & Size & Retained F/G & Flat Combined & Graph Combined \\
\midrule
Original candidate pool & 4 & \pending/\pending & \pending & \pending \\
Add 4 distractors & 8 & \pending/\pending & \pending & \pending \\
Add 16 distractors & 20 & \pending/\pending & \pending & \pending \\
Replace best source, fixed size & 20 & n/a & \pending & \pending \\
\bottomrule
\end{tabular}
\end{table}

\begin{table}[htbp]
\centering
\caption{\textbf{When does executable memory repay its construction cost?}
Measure actual source construction, before/after execution, embedding, and
indexing charges separately from solve-time retrieval and condition execution.
For each benchmark and method, use only first-pass trajectories to report
$[C_{\mathrm{build}}+C_{\mathrm{solve}}(N)]/R(N)$ at the listed cumulative task
counts in the fixed task order. $R(N)$ counts first-pass official resolutions;
solve cost includes target-time construction and failed attempts. Prefixes beyond the manifest,
or with $R(N)=0$, are n/a. For each memory method, report the first observed
prefix with at least the control's resolutions and lower total cost per
resolution, or ``not reached.'' This is the first observed cost crossover,
not a prediction that the advantage persists at every larger prefix.
Use measured costs at these prefixes rather than projecting a success rate.}
\label{tab:planned-amortization}
\small
\begin{tabular}{lccccc}
\toprule
Method & Build USD & USD/solve at 25 & At 50 & At 100 & First crossover \\
\midrule
No memory & n/a & \pending & \pending & \pending & Reference \\
FAISS Flat & \pending & \pending & \pending & \pending & \pending \\
ContextGraph & \pending & \pending & \pending & \pending & \pending \\
\bottomrule
\end{tabular}
\end{table}
\endgroup
